\mode*
\subsection{Classes and objects}

The concept of classes and objects in object-oriented languages is difficult 
and a basic understanding of objects is something that many CS1 students lack 
\autocite{Kaczmarczyk2010}. This is emphasized by \textcite{Ragonis2005OOP} 
who in their study noticed a number of dire misconceptions, two which are 
presented below. 

\begin{enumerate}
    \item An instance of a class can be created within the class' method.

    \item It is possible to define a method which does not access any of the 
class' attributes
\end{enumerate}

Both points concern one object of learning: \emph{what a method is}. The
critical aspect is that a method is simply a function bound to the class; it may
or may not use the instance's attributes, and it may create further instances.
Students who have only seen methods that read and write \mintinline{python}{self}
over-generalise and assume every method must. We separate \enquote{being a
method} from \enquote{using the attributes} by holding the class invariant and
varying only whether the method touches \mintinline{python}{self}.

\begin{description}
  \item[Contrast] The same class, with two equally valid methods: one uses the
    instance's attribute, the other does not. Both are methods; only the use of
    \mintinline{python}{self.name} varies.

    \hfill
    \begin{minipage}[t]{0.45\columnwidth}
      \begin{minted}[highlightlines={5}]{python}
        class Dog:
            def __init__(self, name):
                self.name = name
            def greet(self):
                return f"Woof, I am {self.name}"
      \end{minted}
    \end{minipage}
    \hfill
    \begin{minipage}[t]{0.45\columnwidth}
      \begin{minted}[highlightlines={5}]{python}
        class Dog:
            def __init__(self, name):
                self.name = name
            def sound(self):
                return "Woof"
      \end{minted}
    \end{minipage}
\end{description}


\autocite{Ragonis2005OOP} also found that students have a hard time 
visualising the class as a template for a type of object. Instead the 
students have the 
image of the class as a collection of objects and that the class' methods 
have the power to change, add and delete objects that are class-instances. 
Similar 
misconceptions has been characterised by \textcite{Holland1997}, where the 
students believe that \begin{enumerate*}
    \item an object is a variable that can only hold one or several values of 
the same type and
    \item a class is strictly a data base
\end{enumerate*}. These particular misconceptions can be traced back to the 
first classes the students write, which are often good substitutes to data 
bases 
and therefor shapes these student misconceptions. 

The object of learning is \emph{what a class is}. The critical aspect is that a
class is a \emph{template} (a type) from which independent objects are made ---
not a container that holds its instances, and not a database. A productive
contrast, as the misconception itself suggests, sets a user-defined class against
the built-in types the students already use: \mintinline{python}{str} and
\mintinline{python}{list} are templates, and \mintinline{python}{"hi"} and
\mintinline{python}{[1, 2]} are objects made from them, exactly as
\mintinline{python}{Dog} is a template and \mintinline{python}{Dog("Rex")} is an
object made from it.

\begin{description}
  \item[Contrast] We vary only whether the type is built-in or user-defined,
    holding invariant the relationship \enquote{call the type to make an
    object}. This shows the class is the same kind of thing as the familiar
    built-in types --- a template --- rather than a store of data.

    \hfill
    \begin{minipage}[t]{0.45\columnwidth}
      \begin{minted}{python}
        # str is a template
        s = str("hi")
        t = str("bye")
      \end{minted}
    \end{minipage}
    \hfill
    \begin{minipage}[t]{0.45\columnwidth}
      \begin{minted}{python}
        # Dog is a template
        a = Dog("Rex")
        b = Dog("Fido")
      \end{minted}
    \end{minipage}
    \newline
    The class itself does not hold \mintinline{python}{a} and
    \mintinline{python}{b}: nothing in \mintinline{python}{Dog} records the dogs
    that have been made, just as \mintinline{python}{str} keeps no list of the
    strings in the program. That is what separates a template from a database.

  \item[Generalisation] Keeping the critical value invariant --- a type is a
    template for making objects --- we vary the type across several built-ins
    and the user class, so the principle is seen to hold generally.

    \hfill
    \begin{minipage}[t]{0.3\columnwidth}
      \begin{minted}{python}
        xs = list()
      \end{minted}
    \end{minipage}
    \hfill
    \begin{minipage}[t]{0.3\columnwidth}
      \begin{minted}{python}
        n = int("7")
      \end{minted}
    \end{minipage}
    \hfill
    \begin{minipage}[t]{0.3\columnwidth}
      \begin{minted}{python}
        d = Dog("Rex")
      \end{minted}
    \end{minipage}
\end{description}

The last concept that \textcite{Holland1997} discuss is the concept of 
storing the objects in the programme. Some students believe that the 
attributes of an 
object are the objects identifier, which leads to the misconception that 
there can not be two objects that have the same attributes, and therefor that 
one 
attribute of every object must be unique otherwise the programme will not be 
able to store it. The concept that every object has its own memory space and 
are 
stored separately is hard to grasp for some students \autocite{
Holland1997,Ragonis2005OOP}. 

The object of learning is \emph{object identity and storage}. The critical
aspect is that each object has its own identity and its own memory: two objects
may carry equal attribute values and still be distinct, so an attribute is not an
object's identifier. We hold the attribute values invariant and vary only
whether we have one object or two.

\begin{description}
  \item[Contrast] Two objects with \emph{identical} attributes versus two names
    for the \emph{same} object. On the left \mintinline{python}{a is b} is
    \mintinline{python}{False} and changing one leaves the other untouched ---
    equal attributes do not make one object. On the right both names refer to one
    object, so the change is seen through either.

    \hfill
    \begin{minipage}[t]{0.45\columnwidth}
      \begin{minted}[highlightlines={2}]{python}
        a = Dog("Rex")
        b = Dog("Rex")
        a.name = "Max"
        print(a.name, b.name)
        # Max Rex
      \end{minted}
    \end{minipage}
    \hfill
    \begin{minipage}[t]{0.45\columnwidth}
      \begin{minted}[highlightlines={2}]{python}
        a = Dog("Rex")
        b = a
        a.name = "Max"
        print(a.name, b.name)
        # Max Max
      \end{minted}
    \end{minipage}
    \newline
    The two same-named dogs are stored separately and have different identities,
    which is why a program can hold many objects whose attributes coincide ---
    refuting the belief that one attribute of every object must be unique for it
    to be stored.
\end{description}

